%! Characteristics of Rational Functions — Unit 4, Lesson 4.0 — Homework
\documentclass[10pt]{article}
\usepackage{algebra2-article}
\usepackage{algebra2-boxes}
\newcommand{\TallMath}[1]{$\displaystyle #1\rule[-1.4em]{0pt}{3.2em}$}
\newcommand{\lgrid}[4]{%
  \draw[step=1, linegray!40, very thin] (#1,#3) grid (#2,#4);
  \draw[->, semithick, charcoal] (#1,0)--(#2,0) node[below right, font=\scriptsize]{$x$};
  \draw[->, semithick, charcoal] (0,#3)--(0,#4) node[left, font=\scriptsize]{$y$};}

\begin{document}

\pageheader{Unit 4, Lesson 4.0}{Homework}
\namedateperiod

\vspace{0.10in}

\begin{notesbox}{Practice}
Show how you read each feature. Write every asymptote as an \textbf{equation}.

\begin{enumerate}[leftmargin=1.9em, itemsep=12pt]
  \item \textbf{From the equation only.} Factor where you can, then fill in each row. Write ``none''
        when there is nothing to report.
    \begin{center}
    \renewcommand{\arraystretch}{1.6}
    \begin{tabularx}{\linewidth}{@{}l X X X X@{}}
    \hline
    \textbf{Function} & \textbf{Excluded value(s)} & \textbf{Vert.\ asymptote(s)} & \textbf{Hole} & \textbf{Horiz.\ asymptote} \\
    \hline
    \TallMath{\frac{5}{x+3}} & \blank{1.7cm} & \blank{1.7cm} & \blank{1.7cm} & \blank{1.7cm} \\
    \TallMath{\frac{4x}{x-6}} & \blank{1.7cm} & \blank{1.7cm} & \blank{1.7cm} & \blank{1.7cm} \\
    \TallMath{\frac{x^2-16}{x+4}} & \blank{1.7cm} & \blank{1.7cm} & \blank{1.7cm} & \blank{1.7cm} \\
    \TallMath{\frac{x-2}{x^2-x-2}} & \blank{1.7cm} & \blank{1.7cm} & \blank{1.7cm} & \blank{1.7cm} \\
    \hline
    \end{tabularx}
    \end{center}

  \item \textbf{A graph with two asymptotes.} The graph is $m(x) = \dfrac{x-4}{x-2}$.
    \begin{center}
    \begin{tikzpicture}[scale=0.44]
      \lgrid{-4}{8}{-4}{6}
      \draw[navy, dashed, semithick] (2,-4)--(2,6);
      \draw[navy, dashed, semithick] (-4,1)--(8,1);
      \begin{scope}
        \clip (-4,-4) rectangle (8,6);
        \draw[very thick, forest, domain=-4:1.6, samples=90, smooth]
          plot (\x,{((\x)-4)/((\x)-2)});
        \draw[very thick, forest, domain=2.4:8, samples=90, smooth]
          plot (\x,{((\x)-4)/((\x)-2)});
      \end{scope}
      \fill[charcoal] (4,0) circle (3pt) node[below right, font=\scriptsize]{$(4,0)$};
      \fill[charcoal] (0,2) circle (3pt) node[left, font=\scriptsize]{$(0,2)$};
    \end{tikzpicture}
    \end{center}
    Vertical asymptote: \blank{1.8cm} \quad Horizontal asymptote: \blank{1.8cm} \\[3pt]
    Domain: \blank{4.4cm} \quad Range: \blank{4.4cm} \\[3pt]
    $x$-intercept: \blank{1.6cm} \quad $y$-intercept: \blank{1.6cm} \\[3pt]
    Increasing on: \blank{2.4cm} \; and \; \blank{2.4cm} \\[3pt]
    End behavior: as $x\to\pm\infty$, $m(x)\to\blank{1.0cm}$

  \item \textbf{A graph with a hole.} The graph is $n(x) = \dfrac{x^2-x-6}{x-3}$.
    \begin{center}
    \begin{tikzpicture}[scale=0.46]
      \lgrid{-5}{5}{-3}{7}
      \begin{scope}
        \clip (-5,-3) rectangle (5,7);
        \draw[very thick, forest] (-5,-3)--(5,7);
      \end{scope}
      \draw[forest, very thick, fill=white] (3,5) circle (4.2pt);
    \end{tikzpicture}
    \end{center}
    Factor and simplify: $n(x)=\blank{4.0cm}$, \; for $x\neq\blank{1.0cm}$ \\[3pt]
    The hole is at \blank{1.8cm} \quad Vertical asymptote: \blank{1.6cm} \quad Horizontal asymptote:
    \blank{1.6cm} \\[3pt]
    Domain: \blank{4.4cm} \quad $x$-intercept: \blank{1.6cm} \quad $y$-intercept: \blank{1.6cm}

  \item \textbf{Explain.} A graph may \emph{cross} its horizontal asymptote, but it can \emph{never}
        cross a vertical asymptote. Explain the difference. \writelines{3}
\end{enumerate}
\end{notesbox}

\begin{extensionbox}
\textbf{Reading a model in context.} After a patient takes one dose of a medication, the concentration
in the bloodstream (in mg/L) $t$ hours later is modeled by
\[
  C(t) = \frac{5t}{t^2+1}, \qquad t \ge 0 .
\]
\begin{center}
\renewcommand{\arraystretch}{1.25}
\begin{tabular}{|c|c|c|c|c|c|c|}
\hline
$t$ (hours) & $0$ & $0.5$ & $1$ & $2$ & $4$ & $8$ \\ \hline
$C(t)$ (mg/L) & $0$ & $2$ & $2.5$ & $2$ & $\approx 1.18$ & $\approx 0.62$ \\ \hline
\end{tabular}
\end{center}
\begin{enumerate}[label=(\alph*), leftmargin=1.8em, itemsep=5pt]
  \item Compare the degrees of the numerator and denominator to find the horizontal asymptote.
  \item What does that horizontal asymptote \emph{mean} about the medication over a long time?
  \item From the table, when is the concentration highest, and what is that concentration?
  \item Notice that $C(0)=0$ --- the graph \emph{sits on} its horizontal asymptote at the start. Does
        that contradict anything from today's lesson? Explain.
  \item The denominator $t^2+1$ is \emph{never} zero, so this rational function has \textbf{no}
        vertical asymptote and no excluded values at all. Why, then, is the domain written $t\ge 0$?
\end{enumerate}
\writelines{4}
\end{extensionbox}

\begin{spiralbox}
\textbf{Coming next --- Lesson 4.1: Simplifying Rational Expressions \& Domain Restrictions.} Today
you \emph{read} holes and asymptotes off a graph. Next you will do the algebra that produces them:
factor a rational expression, cancel common factors, and list the restrictions --- always from the
\textbf{original} denominator. That last habit is the one that will save you in Lesson~4.6, when a
``solution'' turns out to be an excluded value in disguise.
\end{spiralbox}

\end{document}
